\documentclass[10pt,a4paper]{amsart}
\usepackage[margin=19mm]{geometry}
\usepackage{amsmath,amssymb,mathtools}
\usepackage[scr=rsfs,cal=euler]{mathalfa}
\usepackage{xfrac}
\usepackage[unicode,hidelinks,bookmarks=false]{hyperref}
\newcommand{\ie}{i.e.\ }
\newcommand{\cf}{cf.\ }
\newcommand{\resp}{respectively\ }
\newcommand{\Subsection}{\subsection}
\providecommand{\nobreakdash}{\nobreak\hbox{-}}
\setlength{\emergencystretch}{2em}
\multlinegap0pt
\usepackage{bm}
\usepackage{bbm}
\usepackage{dsfont}
\usepackage{xspace}
\usepackage{cancel}
\usepackage{mathtools}
\usepackage{amsthm}
\makeatletter
\@ifundefined{theorem}{\newtheorem{theorem}{Theorem}}{}
\@ifundefined{lemma}{\newtheorem{lemma}[theorem]{Lemma}}{}
\makeatother
\title[A study on state spaces in classical Banach spaces]{A study on state spaces in classical Banach spaces}
\author{Soumitra Daptari, Saurabh Dwivedi}
\date{Source: arXiv:2509.12780v1 (CC BY 4.0)}
\begin{document}
\maketitle

\noindent\emph{Source and licence.} A study on state spaces in classical Banach spaces. Version arXiv:2509.12780v1 (CC BY 4.0), distributed under the Creative Commons Attribution 4.0 International licence, as recorded in the arXiv licence metadata for this exact version and shown on the abstract page https://arxiv.org/abs/2509.12780v1. Full rights notice: licenses/arxiv-2509.12780-CC-BY-4.0.txt.\par\smallskip

\noindent\textit{Editorial scope.} Counted unit: the Theorem whose source label is \texttt{Th316} in the article (source0.tex lines 200--202, source label \texttt{Th316}), delivered together with the article's own complete original proof of it (source0.tex lines 203--229, one \texttt{proof} environment of 27 lines). No mathematics is written, repaired or re-derived: statement and proof are the article's own text line for line. Required context retained uncounted: 3t1 (lines 156--158); L1 (lines 193--195). Every definition, display and label of those lines is retained unchanged. Omitted from this package and not redistributed: 1--155, 159--192, 196--199, 230--539. Nothing inside a retained excerpt is omitted or altered: every retained body is delivered exactly as the article's lines are written, so no omission receipt of any kind accompanies this package. Citations: the retained excerpts carry no citation command at all, so no bibliography entry of the article is transcribed and no reference is redistributed. Reference adaptation: none was needed and none was made; every reference inside the delivered unit resolves inside it, so no unresolved reference or citation is produced. Printed slips kept literal: no printed word, symbol, hypothesis or number of the counted statement or of its proof was altered. Layout and class: the article's own class is \texttt{amsart} with packages including amssymb, enumitem, graphicx, fontenc; the wrapper is the lane's pinned \texttt{amsart} editorial wrapper at 10pt on A4, which restates the theorem-like environments the retained text needs and adds the article's own macro definitions that the retained text needs (none), with extra wrapper packages (bm, bbm, dsfont, xspace, cancel, mathtools). The delivered numbering is the unit's own. Assets: no retained span contains a figure environment, a graphics-inclusion command, a PDF-page inclusion, a table or a TikZ picture; no image file is copied into this package, the package declares no asset dependency and no image file is redistributed with it. Licence: version arXiv:2509.12780v1 of this article is distributed under the Creative Commons Attribution 4.0 International licence, as recorded in the arXiv licence metadata for this exact version and shown on the abstract page https://arxiv.org/abs/2509.12780v1. A full rights notice is delivered at licenses/arxiv-2509.12780-CC-BY-4.0.txt. Characters: every retained excerpt is pure ASCII, so the delivered engine, XeLaTeX with its default Latin Modern text font, reports no missing character.
\begin{theorem}\label{3t1}
   Assume $X^*$ is separable. Let $ f \in L^{1}([0,1],X)$ be such that $ \|f\| = 1 $, and $f(t)\neq 0$ almost everywhere. Then $ f $ is smooth in $L^{1}([0,1],X)$ if and only if $ \frac{f(t)}{\|f(t)\|} $ is smooth in $X$ almost everywhere.
\end{theorem}

    \begin{lemma}\label{L1}
        Let $X$ be a separable space and let $f\in L^{1}([0,1],X)$ be a non-zero function such that $\|f\|=1$ and $f(t)\neq0$ almost everywhere. For each $k\in \mathbb{N}$, the set $A_{k}=\{t\in [0,1]:d^*(x_{t}^*, y_{t}^*)\geq \frac{1}{k}\textnormal{ for some }x_{t}^*, y_{t}^*\in S_{\frac{f(t)}{\|f(t)\|}}\}$ is a measurable set.
    \end{lemma}

\begin{theorem}\label{Th316}
    Assume $X^*$ is separable. Let $f\in L^{1}([0,1],X)$ be such that $\|f\|=1$, and $f(t)\neq0$ almost everywhere. Then $S_f$ is norm compact if and only if $f$ is smooth.
\end{theorem}

\begin{proof}
 Since Lebesgue $\sigma$-field is complete, by discarding null sets if necessary, we consider all the functions are everywhere defined. Since $f\neq0$ almost everywhere, we will assume $f(t)\neq0$ for all $t\in [0,1]$. If $f$ is a smooth point, then $S_{f}$ is a singleton, hence a norm compact set. If $f$ is not smooth, then we show that $S_{f}$ is not norm compact by constructing a sequence in $S_{f}$ with no convergent subsequence. Fix $h\in S_{f}$, and let $A=\{t\in[0,1]:{\frac{f(t)}{\|f(t)\|}}\text{ is not smooth}\}$. From Theorem~\ref{3t1}, we have $\mu(A)>0$. For $k\in \mathbb{N}$, we define $A_{k}=\{t\in [0,1]:d^*(x_{t}^*, y_{t}^*)\geq \frac{1}{k}\textnormal{ for some }x_{t}^*, y_{t}^*\in S_{\frac{f(t)}{\|f(t)\|}}\}$. From Lemma~\ref{L1}, we have that each $A_{k}\subseteq [0,1]$ is measurable. It is easy to see that $A=\bigcup_{k\in \mathbb{N}}A_{k}$. Since $\mu(A)>0$, we have that $\mu(A_{k})>0$ for some $k\in \mathbb{N}$. Let $Z=A_{k}$, $2\epsilon_{0}=\frac{1}{k}$ and $B_{d^*}^{o}(h(t), \epsilon_{0})$ denotes the open $\epsilon_{0}$-ball centred at $h(t)$ in the weak$^*$ metric. Similarly, $B_{\|.\|}^{o}(h(t), \epsilon_{0})$ denotes the open $\epsilon_{0}$-ball centred at $h(t)$ in norm. For each $t\in Z$, the set $S_{\frac{f(t)}{\|f(t)\|}}\setminus B_{d^*}^{o}(h(t), \epsilon_{0})$ is non-empty and measurable. Since the metric $d^*$ satisfies $d^*(x^*,y^*)\leq \|x^*-y^*\|$ for all $x^*,y^*\in X_{1}^* $, we must have that the set $S_{\frac{f(t)}{\|f(t)\|}}\setminus B_{\|.\|}^{o}(h(t), \epsilon_{0})$ is non-empty and measurable for each $t\in Z$. Since $\mu$ is non-atomic and $\mu(Z)>0$, we can choose a sequence $\{Z_{n}\}_{n\in \mathbb{N}}$ of disjoint subsets of $Z$ such that $\mu(Z_{n})>0$ for each $n\in \mathbb{N}$. As $X^*$ is a separable space, we have that $S(X^*)$ is a Polish space with respect to the norm topology.

Define $ F_1: [0,1] \to 2^{S(X^*)} $ by
\[
F_{1}(t) =
\begin{cases} 
S_{\frac{f(t)}{\|f(t)\|}}\setminus B_{\|.\|}^{o}(h(t), \epsilon_{0}), & t \in Z_{1} , \\
S_{\frac{f(t)}{\|f(t)\|}}, & t \notin Z_{1}.
\end{cases}
\]
Then
\[
G(F_1) =\{(t,x^*):x^*(f(t))=\|f(t)\|\} \setminus \{(t,x^*):\|x^*-h(t)\|<\epsilon_{0}, t \in Z_1\}.
\]

We already have shown in the proof of Theorem~\ref{3t1} that the set $\{(t,x^*):x^*(f(t))=\|f(t)\|\}$ is a Borel set. Since the mappings $(t,x^*)\rightarrow x^*-h(t)$ and $\ (x^*-h(t))\rightarrow\|x^*-h(t)\|$ are measurable and continuous, respectively, we have $\{(t,x^*):\|x^*-h(t)\|<\epsilon_{0}, t \in Z_1\}$ is a measurable set. Hence $ G(F_1) $ is a measurable set. Applying Von Neumann's selection theorem to $G(F_{1})$ as in Theorem~\ref{3t1}, we get a selection $h^{'}_1: [0,1] \to S(X^*)$ such that $\|h^{'}_1(t)- h(t)\|\geq\epsilon_{0}\text{ for all  } t \in Z_1$. This implies $ \|h^{'}_1- h\|\geq \epsilon_{0}$. Define $ h_1: [0,1] \to X^* $ by
\[
h_{1}(t) =
\begin{cases} 
h^{'}_{1}(t), & t \in Z_{1} , \\
h(t), & t \notin Z_{1}.
\end{cases}
\]
Fix $h_{0}=h$. Then it is easy to see that $h_{0}$ and  $h_1$ are two distinct elements in $ S_f $ such that $\|h_1- h_{0}\|\geq\epsilon_{0}$ and $h_{0}(t)=h_{1}(t) \text{ for all } t\in Z_{1}^{c}$.
Continuing this process for each $Z_{n}$, we get a sequence $\{h_n\}_{n\geq0}$ such that $\|h_n-h_m\|\geq\epsilon_{0}$ for all $n\neq m$. Hence $S_{f}$ is not norm compact. This completes the proof.
\end{proof}

\end{document}
